\documentclass[output=paper]{langscibook}
\ChapterDOI{10.5281/zenodo.22880246}
\title{\textit{Koverau, kururu, koopu} -- Three kinds of bamboo}
\author{Ruth Sima Rigamu and Joyce Maion}
\abstract{}
\begin{document}
\maketitle\label{chap:10}
\begin{figure}[H]
    \captionsetup[subfigure]{oneside,margin={0cm,0cm}}
    \subfloat[]{\includegraphics[trim=0 0 0 0,clip,width=.33\linewidth]{figures/10_1.jpg}\label{fig:10_1}}
    \subfloat[]{\includegraphics[trim=0 0 0 50,clip,width=.34\linewidth]{figures/10_2.jpg}\label{fig:10_2}}
    \subfloat[]{\includegraphics[trim=0 0 0 0,clip,width=.33\linewidth]{figures/10_3.jpg}\label{fig:10_3}}
    \caption{(a) \textit{Koverau}, (b) \textit{Kururu}, (c) \textit{Koopu}}
\end{figure}
\begin{figure}[H]
    \includegraphics[width=\linewidth]{figures/10_4.png}
    \caption{\textit{Koverau}}
    \label{fig:10_4}
\end{figure}
\begin{figure}[H]
    \includegraphics[width=\linewidth]{figures/10_5.jpg}
    \caption{\textit{Makaree} -- Fishing spear called \textit{makaree}}
    \label{fig:10_5}
\end{figure}
\begin{xltabular}{\textwidth}{lXX}
    (\ref{ex:text6-1})&\textit{O koverau na pura nana rasuu.}&The \textit{koverau} bamboo grows in the jungle.\\
    (\ref{ex:text6-2})&\textit{O }koverau\textit{ saka tau pura vatotobin tamuana bata haana.}&The \textit{koverau} bamboo might not always grow straight.\\
    (\ref{ex:text6-3})&\textit{Tau hoa vanaavihu batana bara tau obete gunaha suku batana bona maa hum gunaha.}&It might grow in the opposite direction and lie on the ground following the slope of a mountain downward.\\
    (\ref{ex:text6-4})&\textit{Ei be pura tea maa vaereere.}&This means that it grows on the edges of mountains.\\
    (\ref{ex:text6-5})&\textit{Pura vageegeno batari amaa meha.}&Some grow crooked.\\
    (\ref{ex:text6-6})&\textit{O koverau na vaarutaa ki nana bono kururu.}&The \textit{koverau} bamboo is smaller than the \textit{kururu} bamboo.\\
    (\ref{ex:text6-7})&\textit{O naono o mataa, na paku raara bono tahitahi.}&It is a good plant, we make spears (from it).\\
    (\ref{ex:text6-8})&\textit{Bea otei tau nao tea mosi koverau,}&When a man is about to cut \textit{koverau} bamboos,\\
    (\ref{ex:text6-9})&\textit{toro tara vamataa kurus maa bona maa aba vai}&he must have a very good look at the ones\\
    (\ref{ex:text6-10})&\textit{to pasi mosi maa nae mene o makaree teve.}&that he will cut for his fishing spear.\\
    (\ref{ex:text6-11})&\textit{Be dee voosu maa bona dee koverau teve,}&When he brings his load of \textit{koverau} bamboos home,\\
    (\ref{ex:text6-12})&\textit{ore paa vaasun bona matana inu teori,}&he stands them up in front of their house,\\
    (\ref{ex:text6-13})&\textit{ei ore roava nana.}&so that they are exposed to the sun.\\
    (\ref{ex:text6-14})&\textit{Be mahaka, ore paa aeae, ore paa nana mana bari teo suraa.}&When they dry up, they become light, and then they (the men) warm them over the fire.\\
    (\ref{ex:text6-15})&\textit{Pasi nana bata nom ori, beori menamenaa vaatotobin batari bari.}&They will warm them, while they are bending them straight.\\
    (\ref{ex:text6-16})&\textit{Beve totobin, ore paa kahu bono kosi no koverau,}&When they (the bamboos) are straight, they scrape the bottom of the \textit{koverau} bamboos,\\
    (\ref{ex:text6-17})&\textit{vaamatatopo vo tea vaatopo bono kuna sikiri.}&and prepare them for sticking the spikes (into them).\\
    (\ref{ex:text6-18})&\textit{Vaatopo vai bona maa kuna sikiri,}&Now they stick in the spikes,\\
    (\ref{ex:text6-19})&\textit{ore paa pam bari bona kanono.}&and fasten them with a string.\\
    (\ref{ex:text6-20})&\textit{Be kavara, erau are paa pokaa e bono makaree.}&When this is finished, then we call it \textit{makaree} `fishing spear'.\\
    (\ref{ex:text6-21})&\textit{Bea abana nao tahii, toro ato momohu tea makaree teori mene tahi vamate bona iana vai a beera to pasi rahi riori tea kanono.}&When the men go fishing, they must first get their fishing spears for spearing a big fish to death that they will pull (out of the water) with a string.\\  
\end{xltabular}
\vspace{-.5cm}
\begin{figure}[H]
    \includegraphics[width=.8\linewidth]{figures/10_6.jpg}
    \caption{\textit{Ato teo makaree} -- Holding the \textit{makaree} spear}
    \label{fig:10_6}
\end{figure}
\begin{figure}[H]
    \includegraphics[width=.8\linewidth]{figures/10_7.png}
    \caption{\textit{Kururu}}
    \label{fig:10_7}
\end{figure}
\begin{xltabular}{\textwidth}{lXX}
    (\ref{ex:text6-22})&\textit{O kururu na vaapeha mi nana bono koverau.}&The \textit{kururu} bamboo is similar to the \textit{koverau} bamboo.\\
    (\ref{ex:text6-23})&\textit{Amaa naono vai to tei me rori bona maa navanava riori: koverau, kururu, koopu, tovu, tabuana.}&The plants that have bars (separated by knots) are: the \textit{koverau}, \textit{kururu} and \textit{koopu} bamboos, the \textit{tovu} sugar cane and the wild sugar cane.\\
    (\ref{ex:text6-24})&\textit{O kururu na isuisuvu ruene ni riori mee.}&The \textit{kururu} bamboo is used for fetching water.\\
    (\ref{ex:text6-25})&\textit{Beori tau isuvu ruene ni bari, na dona rori bona kukan a nava kururu.}&When they are going to fetch water with it, they pierce (the knots between) three bars.\\
    (\ref{ex:text6-26})&\textit{Murinae, eori repaa gono bene ta naono eta bebeahu bara kikis}&For this they get a long and strong stick\\
    (\ref{ex:text6-27})&\textit{to tahitahi nibusu nao riori bona buaku a nava kururu vaa gina,}&with which they pierce holes into the knots of the two top \textit{kururu} bars,\\
    (\ref{ex:text6-28})&\textit{beori goe bona nava vaa taeno kururu.}&while leaving the bar at the bottom of the \textit{kururu} bamboo intact.\\
    (\ref{ex:text6-29})&\textit{A ruene re goe tea oro tavusu.}&Then the water would not flow out.\\    
    (\ref{ex:text6-30})&\textit{O koopu saka pura vamataa haana rapana kasuana ge tea suiri,}&The \textit{koopu} bamboo does not grow well near the beach or in the lowlands,\\
    (\ref{ex:text6-31})&\textit{na pura vamataa nana tea voora.}&it grows well in the mountains.\\
    (\ref{ex:text6-32})&\textit{Meha maa vuri a ta vaan vai to tei rori tea maa vaan vahiava na vaamahaka rori bona tabaan teori komano koopu.}&Sometimes the people who live in the mountain villages cook their food in the \textit{koopu} bamboo.\\
    (\ref{ex:text6-33})&\textit{Na mosi maari bona maa nava koopu,}&They cut some bars of \textit{koopu} bamboo,\\
\end{xltabular}
\begin{figure}[H]
    \captionsetup[subfigure]{oneside,margin={0cm,0cm}}
    \subfloat[]{\includegraphics[trim=0 30 0 26,clip,width=.51\linewidth]{figures/10_8.jpg}\label{fig:10_8}}
    \subfloat[]{\includegraphics[trim=0 30 0 0,clip,width=.49\linewidth]{figures/10_9.jpg}\label{fig:10_9}}
    \caption{(a) \textit{Koopu}, (b) \textit{Taboo a inu} -- Walling a house}
\end{figure}
\begin{xltabular}{\textwidth}{lXX}
    (\ref{ex:text6-34})&\textit{eori repaa apeo bona tabaan tabae to rake riori tea apeo: e guu, manii, kaukau, gea tabae to rake riori tea an.}&and then they cook in the bamboo whatever food they like to cook in the bamboo: pork, possum, sweet potatoes or what they like to eat.\\
    (\ref{ex:text6-35})&\textit{O koopu vai o peree to apeo riori bona tabaan.}&It is the green \textit{koopu} bamboo in which they cook the food.\\
    (\ref{ex:text6-36})&\textit{Eori repaa vaamahaka bari teo suraa.}&Then they let it simmer on the fire.\\
    (\ref{ex:text6-37})&\textit{Beori vaatoka bono koopu teo suraa,}&When they put the \textit{koopu} bamboo onto the fire,\\
    (\ref{ex:text6-38})&\textit{toro vavereverete bata ni rori bono koopu to mene asi.}&they must repeatedly turn the \textit{koopu} bamboo around so that it won't burn.\\
    (\ref{ex:text6-39})&\textit{Eara na voni rara a koopu tea ta vahiava mene paku bono takinima na maa inu teara.}&We buy the \textit{koopu} bamboo from the mountain people to make the walls of our houses.\\
    (\ref{ex:text6-40})&\textit{Beara taboo a inu teara bono takinima ore paa tara vamataa.}&When we wall our house with bamboo walls,  it looks beautiful.\footnotemark\\
\end{xltabular}
\footnotetext{See Chapter \ref{chap:19}.}
\end{document}